\documentclass[conference]{IEEEtran}

\IEEEoverridecommandlockouts
% The preceding line is only needed to identify funding in the first footnote. If that is unneeded, please comment it out.
\usepackage{cite}
\usepackage{amsmath,amssymb,amsfonts}
% \usepackage{algorithm}
\usepackage[ruled,vlined]{algorithm2e}
% \usepackage{algorithmic}
\usepackage{algpseudocode}
\usepackage{graphicx}
\usepackage{textcomp}
\usepackage{xcolor}
\usepackage{booktabs}
\usepackage{array}
\usepackage{makecell}
\usepackage{multirow}
\usepackage{subfig}
\def\BibTeX{{\rm B\kern-.05em{\sc i\kern-.025em b}\kern-.08em
    T\kern-.1667em\lower.7ex\hbox{E}\kern-.125emX}}
\begin{document}

\title{SLNN-Assisted Decoding of a 7/9-Rate Sparse Code for Reliable STT-MRAM\\}
% \title{Single Label Neural Network Assisted Reliability Enhancement for 7/9-rate Sparse in STTMRAM\\}
% {\footnotesize \textsuperscript{*}Note: Sub-titles are not captured in Xplore and
% should not be used}
% \thanks{Identify applicable funding agency here. If none, delete this.}
% }

\author{
\IEEEauthorblockN{
Nhon Q. Nguyen,
Huy V. Giang,
Thanh-Loc Nguyen-Van,
Tan Do-Duy
}

\IEEEauthorblockA{Ho Chi Minh City University of Technology and Engineering (HCM-UTE), Vietnam
}

}

\maketitle

\begin{abstract}
This paper proposes a concatenated sparse coding and BCH coding framework to improve the reliability of Spin-Transfer Torque Magnetic Random Access Memory (STT-MRAM). The framework employs the existing 7/9-rate sparse code to exploit the asymmetric write-error characteristic of STT-MRAM by reducing the number of $(0 \rightarrow 1)$ write transitions, while a BCH code corrects the residual memory errors. To overcome the limitation of conventional maximum-likelihood (ML) sparse decoding, a Single-Label Neural Network (SLNN)-based soft-decision decoder is developed to exploit the reliability information of the received signals. By improving sparse decoding accuracy, the proposed decoder reduces error propagation to the BCH decoder and enhances the overall decoding performance. Simulation results demonstrate that the proposed SLNN-based decoder outperforms the conventional ML decoder. Furthermore, the proposed concatenated coding framework achieves lower bit error rate (BER) than standalone BCH coding and conventional sparse coding.
\end{abstract}


\begin{IEEEkeywords}
STT-MRAM, concatenated codes, sparse codes, BCH codes, neural network decoding, soft-decision decoding.
\end{IEEEkeywords}



\section{Introduction}
Spin-Transfer Torque Magnetic Random Access Memory (STT-MRAM) is a promising non-volatile memory technology because of its high density, fast access, excellent endurance, and compatibility with emerging Internet of Things (IoT) and artificial intelligence (AI) applications \cite{girard2020survey}. However, ensuring reliable data storage remains a major challenge as write- and read-induced failures become increasingly pronounced with continued technology scaling \cite{hemaram2025asymmetric}.

These reliability issues mainly originate from the stochastic switching behavior of magnetic tunnel junctions (MTJs). Process variations in MTJ geometry and resistance alter the switching current distribution, resulting in significant cell-to-cell variability. Since practical memory systems employ a fixed write pulse width, some MTJs fail to complete switching within the allocated write duration, causing write failures. Thermal fluctuations, read disturbance, and manufacturing imperfections further aggravate these effects, making reliability enhancement increasingly difficult as device dimensions continue to shrink \cite{wu2020survey}.

To combat these reliability degradations, channel coding techniques have been widely adopted for memory protection. Existing approaches can be broadly classified into constrained coding, error-correcting codes (ECCs), and their combination \cite{nguyen2026error}. Constrained coding reshapes stored data to suppress error-prone patterns, whereas ECCs, such as Hamming and BCH codes, detect and correct memory errors. Nevertheless, conventional ECCs provide symmetric protection for both bit-flipping directions and therefore cannot fully exploit the asymmetric write-error characteristic of STT-MRAM, where $(0 \rightarrow 1)$ transitions are considerably more vulnerable than $(1 \rightarrow 0)$ transitions.

To exploit this asymmetry, sparse coding has been proposed as an effective constrained coding technique. By reducing the number of stored ``1'' bits, sparse coding decreases the occurrence of $(0 \rightarrow 1)$ write transitions and consequently lowers the write failure probability \cite{nguyen2021design}. In particular, the 7/9-rate sparse code significantly improves write reliability. However, sparse coding does not provide error-correction capability, and its performance degrades when multiple bit errors occur. Therefore, recent studies have investigated concatenating sparse coding with BCH codes, where sparse coding suppresses asymmetric write failures while BCH coding corrects the remaining memory errors \cite{nguyen2026error}. Despite this complementary protection, the sparse decoder is still typically implemented using conventional maximum-likelihood (ML) decoding. Because ML decoding relies solely on hard decisions, it cannot exploit the reliability information available from the read signals, allowing sparse decoding errors to propagate to the BCH decoder and degrade the overall decoding performance.

Motivated by this limitation, this paper proposes a Single-Label Neural Network (SLNN)-based decoder for the 7/9-rate sparse code and integrates it with BCH coding in a unified reliability enhancement framework. Unlike the conventional ML decoder, the proposed SLNN-based decoder exploits soft information from the received signals to improve sparse decoding accuracy. Combined with BCH error correction, the proposed framework effectively mitigates both asymmetric write failures and residual memory errors, thereby significantly enhancing the overall reliability of STT-MRAM.

The main contributions of this paper are summarized as follows.
%
\begin{itemize}
    \item We propose an SLNN-based soft-decision decoder for the existing 7/9-rate sparse code. By exploiting the reliability information of the received signals, the proposed decoder achieves higher decoding accuracy than the conventional maximum-likelihood (ML) decoder.

    \item We develop a concatenated sparse coding and BCH coding framework for STT-MRAM reliability enhancement, where sparse coding suppresses asymmetric write failures and BCH coding corrects the residual memory errors.

    \item Simulation results show that the proposed SLNN-based decoder outperforms the conventional ML decoder, while the proposed concatenated coding framework achieves lower BER than standalone BCH coding and conventional sparse coding.
\end{itemize}

%% The remainder of this paper is organized as follows. 


\section{Channel model for STT-MRAM}
\subsection{STT-MRAM storage technology}

As illustrated in Fig.~\ref{figcell}, an STT-MRAM cell employs a conventional 1T--1MTJ architecture, consisting of one magnetic tunnel junction (MTJ) and one access nMOS transistor. The MTJ functions as the non-volatile storage element, while the access transistor controls both read and write operations \cite{cai2017cascaded}. The MTJ comprises two ferromagnetic layers separated by a thin insulating tunnel barrier. The reference layer (RL) has a fixed magnetization direction, whereas the free layer (FL) can switch its magnetization under the spin-transfer torque effect. The relative alignment of these two layers determines the resistance of the MTJ. When the RL and FL are aligned in the parallel (P) state, the MTJ exhibits a low resistance, denoted by $R_P$, which represents logic 0'. In contrast, when they are in the antiparallel (AP) state, the MTJ has a high resistance, denoted by $R_{AP}$, corresponding to logic 1'. Therefore, binary information is stored by distinguishing between these two resistance states \cite{nguyen2023improving}.

\begin{figure}[t]
\centerline{\includegraphics [scale=0.12]{fig/01.drawio.pdf}}
\caption{STT-MRAM structure. (a) STT-MRAM cell. (b) Parallel (P) state. (c) Antiparallel (AP) state \cite{nguyen2023improving}.}
\label{figcell}
\end{figure}



\begin{figure}[t]
\centerline{\includegraphics [scale=0.12]{fig/03.drawio.pdf}}
\caption{Writing process in STT-MRAM. (a) Bit-0 writing process. (b) Bit-1 writing process \cite{nguyen2021design}.}
\label{figwriteread}
\end{figure}

The write operation in STT-MRAM is inherently asymmetric, as illustrated in Fig.~\ref{figwriteread}. To write a logic `0', $V_{DD}$ is applied to the WL and BL while the source line (SL) is grounded, driving the MTJ to the low-resistance parallel (P) state. To write a logic `1', the current direction is reversed, switching the MTJ to the high-resistance antiparallel (AP) state. Because of the asymmetric switching behavior of the MTJ, the transition to the AP state generally requires more switching energy and a longer switching time than the transition to the P state. Consequently, the write error probability of bit-1 can be up to three orders of magnitude higher than that of bit-0, making write reliability a major challenge in STT-MRAM-based memory systems \cite{nguyen2021design}.

During the read operation, the word line (WL) is first activated, and a sensing current, much smaller than the write current, is passed through the MTJ to avoid disturbing the stored data. The resulting resistance is then compared with a reference threshold to determine the stored bit value. Read reliability mainly depends on the tunneling magnetoresistance (TMR) ratio, which measures the resistance difference between the P and AP states. A higher TMR ratio provides a larger sensing margin, thereby reducing read decision errors and improving the overall reliability of the memory system \cite{nguyen2021design}.




\subsection{Cascaded channel model}

To evaluate signal processing techniques for STT-MRAM, an accurate channel model is required to characterize the physical error mechanisms of the memory. Several channel models have been proposed in the literature \cite{cai2017cascaded,li2016efficient}. In this work, the cascaded STT-MRAM channel model introduced by Cai and Immink \cite{cai2017cascaded} is adopted.

\begin{figure}[!t]
\centering
\includegraphics[width=\columnwidth]{fig/04_v2.drawio.pdf}
\caption{Cascaded STT-MRAM channel model under the write-1 operation \cite{cai2017cascaded}.}
\label{fig3}
\end{figure}

Let $P_1$, $P_0$, and $P_r$ denote the write error probability for $1\rightarrow0$ transitions, the write error probability for $0\rightarrow1$ transitions, and the read disturb error rate, respectively. As shown in Fig.~\ref{fig3}, the channel consists of three cascaded components. A binary asymmetric channel (BAC) models asymmetric write errors with crossover probabilities $P_0/2$ and $P_1/2$. The read disturb error, which is unidirectional, is modeled by a Z channel. Finally, read decision errors are represented by a Gaussian mixture channel (GMC), where the MTJ resistances are modeled as Gaussian random variables: $R_0 \sim \mathcal{N}(\mu_0,\sigma_0)$ for the low-resistance state and $R_1 \sim \mathcal{N}(\mu_1,\sigma_1)$ for the high-resistance state.

For analysis, the BAC and Z channels are combined into an equivalent BAC-ZC channel with crossover probabilities

\begin{equation}
\begin{aligned}
p_0 &= \frac{P_0}{2} + \left(1 - \frac{P_0}{2}\right) P_r, \quad
q_0 = \left(1 - \frac{P_0}{2}\right)(1 - P_r), \\
p_1 &= \frac{P_1}{2} + (1 - P_r), \quad
q_1 = \left(1 - \frac{P_1}{2}\right) + \frac{P_1}{2} P_r.
\end{aligned}
\end{equation}

where $p_0$ and $p_1$ are the crossover probabilities for $0\rightarrow1$ and $1\rightarrow0$ transitions, respectively, while $q_0$ and $q_1$ denote the probabilities of correct transitions ($0\rightarrow0$ and $1\rightarrow1$).




\section{Proposed SLNN-Assisted Reliability Enhancement Framework}
\subsection{System model}

Fig.~\ref{figsystem} illustrates the block diagram of the proposed system. The framework combines an outer BCH code, which provides error correction capability, with an inner 7/9-rate sparse code that shapes the input data to reduce write errors in STT-MRAM.

At the transmitter, the user data $\mathbf{u}$ is first encoded by the BCH encoder to generate the codeword $\mathbf{c}_1$. The BCH code provides error-correction capability against channel errors. The encoded sequence is then processed by the 7/9-rate sparse encoder to produce the sparse codeword $\mathbf{c}_2$. By reshaping the input data into sparse patterns, the sparse encoder reduces the number of high-energy write operations, thereby lowering the write error probability in STT-MRAM.

The sparse codeword is transmitted through the cascaded STT-MRAM channel, which models asymmetric write errors, read disturbance errors, and read decision errors. At the receiver, the received signal $\mathbf{r}$ is first decoded by the proposed SLNN-based sparse decoder. Unlike conventional hard-decision decoding, the proposed decoder exploits the soft information of the received signal to accurately recover the sparse codeword, producing the estimated BCH codeword $\hat{\mathbf{c}}_1$.

Finally, $\hat{\mathbf{c}}_1$ is passed to the BCH decoder, which corrects the remaining residual errors to recover the original user data $\hat{\mathbf{u}}$. By combining data shaping through the 7/9-rate sparse code, neural-assisted sparse decoding, and BCH error correction, the proposed framework effectively improves the reliability of STT-MRAM. For brevity, the detailed encoding and decoding procedures of the BCH and 7/9-rate sparse codes are omitted in this paper, as they have been comprehensively described in \cite{lin2001error} and \cite{nguyen2021design}, respectively.




\begin{figure*}[t]
\centerline{\includegraphics [scale=0.15]{fig/figsys_02.drawio.pdf}}
\caption{Block diagram of the proposed system.}
\label{figsystem}
\end{figure*}


\subsection{Single-label neural network-based 7/9-rate sparse decoder}



% \begin{table}[tb]
% \centering
% \caption{Summary of the considered BCH–sparse concatenated coding schemes and their corresponding overall coding rates.}
% \label{tab4}
% \renewcommand{\arraystretch}{1.2}
% \begin{tabular}{|c|c|c|}
% \hline
% No. 
% & Coding scheme 
% & Overall rate 
%  \\
% \hline\hline

% 1 & BCH(127,106,3) + 7/9-rate sparse code
% & 0.649 
%  \\ \hline

% 2 & BCH(127,99,4) + 7/9-rate sparse code
% & 0.606
%  \\ \hline

% 3 & BCH(127,106,3) + 4/9-rate sparse code
% & 0.371
%  \\ \hline

% 4 & BCH(127,99,4) + 4/9-rate sparse code 
% & 0.346
%  \\
% \hline
% \end{tabular}
% \end{table}



\section{SIMULATION RESULTS AND DISCUSSION}




\subsection{Simulation setup}
\subsection{Simulation results}


\begin{figure}[t]
\centerline{\includegraphics [scale=0.6]{fig/BER_vs_sigma.png}}
\caption{BER performance versus the coefficient of variation $(\sigma_0/\mu_0)$.}
\label{ber_vs_sigma}
\end{figure}


\begin{figure}[t]
\centerline{\includegraphics [scale=0.6]{fig/FER_vs_sigma.png}}
\caption{FER performance versus the coefficient of variation $(\sigma_0/\mu_0)$.}
\label{fer_vs_sigma}
\end{figure}



\section{Conclusion}

This paper 

\section*{Acknowledgment}







% -------------------------------------------------------------------------
\bibliographystyle{IEEEtran} % ieeetr
\bibliography{refs}





\vspace{12pt}
% \color{red}
% IEEE conference templates contain guidance text for composing and formatting conference papers. Please ensure that all template text is removed from your conference paper prior to submission to the conference. Failure to remove the template text from your paper may result in your paper not being published.

\end{document}
